\documentclass[11pt]{article}

\usepackage[T1]{fontenc}
\usepackage[utf8]{inputenc}
\usepackage{amsmath,amssymb,mathtools,mathrsfs}
\usepackage[margin=0.9in]{geometry}
\usepackage{microtype}
\usepackage{needspace}

\setlength{\parindent}{0pt}
\setlength{\parskip}{0.58em}

\begin{document}

\section*{How the Proof Works}

A vortex tube lengthens when the surrounding strain pulls it along its axis. Incompressibility prevents that lengthening from creating volume, so the cross-section contracts and the rotating core narrows. The fluid is forcing its velocity to turn through a smaller region.

That sharpening takes place inside the three-dimensional Navier--Stokes equation:
\[
\partial_tu+(u\cdot\nabla)u+\nabla p=\nu\Delta u,
\qquad \nabla\cdot u=0,
\qquad u(\cdot,0)=u_0.
\]
The nonlinear motion stretches and concentrates the tube; viscosity diffuses its spatial variation. Curling the equation shows both actions on the vorticity \(\omega=\nabla\times u\):
\[
\partial_t\omega+(u\cdot\nabla)\omega
=(\omega\cdot\nabla)u+\nu\Delta\omega.
\]
Where the vortex is aligned with an amplifying direction of the strain,
\[
\omega\cdot((\omega\cdot\nabla)u)>0,
\]
the stretching part raises \(\lvert\omega\rvert\) while the core narrows. Since vorticity is one spatial derivative of velocity, the velocity changes across less distance and its gradients steepen. The global regularity problem asks whether this sharpening can keep concentrating a smooth velocity field until it becomes singular in finite time.

Viscosity acts more strongly on a narrower core, but so does the nonlinear term. Scaling measures the two changes together. For \(\lambda>1\), compress the entire solution by setting
\[
u_\lambda(x,t)=\lambda u(\lambda x,\lambda^2t),
\qquad
p_\lambda(x,t)=\lambda^2p(\lambda x,\lambda^2t).
\]
A structure of width \(r\), velocity scale \(U\), and duration \(\tau\) becomes one of width \(r/\lambda\), velocity scale \(\lambda U\), and duration \(\tau/\lambda^2\). The nonlinear and viscous terms change as
\[
\frac{U^2}{r}
\longmapsto
\lambda^3\frac{U^2}{r},
\qquad
\nu\frac{U}{r^2}
\longmapsto
\lambda^3\nu\frac{U}{r^2}.
\]
Every term in the equation acquires the factor \(\lambda^3\). Nonlinear sharpening and viscous diffusion strengthen in exact proportion. Shrinking the scale gives neither one an advantage.

Kinetic energy misses this concentration. The compressed core occupies \(\lambda^{-3}\) of the volume while the square of its velocity contributes \(\lambda^2\), so
\[
\|u_\lambda(\cdot,t)\|_2^2
=\lambda^{-1}\|u(\cdot,\lambda^2t)\|_2^2.
\]
The gradients can grow while the core's contribution to the energy falls. Across the homogeneous Sobolev scale,
\[
\|u_\lambda(\cdot,t)\|_{\dot H^s}
=\lambda^{s-\frac12}
\|u(\cdot,\lambda^2t)\|_{\dot H^s}.
\]
Concentration makes the norm smaller below \(s=\tfrac12\) and larger above it. At \(s=\tfrac12\), it does neither. This is the critical height: narrowing the flow gives this measurement no automatic gain and no automatic loss. Write
\[
E_s(t):=\frac12\|\Lambda^su(t)\|_2^2,
\qquad
H(t):=E_{1/2}(t).
\]

The compressed clock leaves room for a finite-time cascade. If each stage contracts by a fixed factor \(0<\rho<1\), then
\[
r_j=\rho^jr_0,
\qquad
U_j=\rho^{-j}U_0,
\qquad
\tau_j=\rho^{2j}\tau_0,
\]
and
\[
\sum_{j=0}^{\infty}\tau_j
=\tau_0\sum_{j=0}^{\infty}\rho^{2j}
=\frac{\tau_0}{1-\rho^2}<\infty.
\]
Infinitely many narrower and faster stages fit inside a finite interval. This is the Zeno cascade. It does not construct a singular solution. It shows what finite energy and the scale balance do not exclude: concentration can intensify while its successive time scales have a finite sum.

Suppose, for contradiction, that the maximal smooth solution ends at \(T_*<\infty\).

In order for a singularity to form, velocity has to blow up. In order for that to happen, its acceleration would also have to blow up, which means its jerk has to blow up, and so on, and so on:
\[
u,\qquad \partial_tu,\qquad \partial_t^2u,\qquad \ldots
\]
What you would expect to see in this pattern, if a singularity were to form, is less spatial regularity as you go up the time derivatives, not more. Navier--Stokes ties each temporal step to a spatial one.

The first three steps are
\[
\partial_tu
=\nu\Delta u-(u\cdot\nabla)u-\nabla p,
\]
\[
\partial_t^2u
=\nu\Delta\partial_tu
-(\partial_tu\cdot\nabla)u
-(u\cdot\nabla)\partial_tu
-\nabla\partial_tp,
\]
\[
\begin{aligned}
\partial_t^3u={}&\nu\Delta\partial_t^2u
-(\partial_t^2u\cdot\nabla)u
-2(\partial_tu\cdot\nabla)\partial_tu \\
&-(u\cdot\nabla)\partial_t^2u
-\nabla\partial_t^2p.
\end{aligned}
\]
The complete pattern is
\[
\partial_t^{\,n+1}u
=\nu\Delta\partial_t^nu
-\sum_{a=0}^{n}\binom na
\bigl(\partial_t^au\cdot\nabla\bigr)\partial_t^{\,n-a}u
-\nabla\partial_t^np.
\]
Only now set
\[
V_n:=\partial_t^nu,
\qquad
Q_n:=\partial_t^np.
\]
Every rise in temporal order brings a Laplacian of the preceding response. Viscosity links the time ladder to spatial derivatives. The nonlinear products spread across the lower temporal orders, and pressure differentiates with them, so every row remains the complete equation.

The time ladder and the critical height meet in the exact balance
\[
E_s'=-2\nu E_{s+1}+\mathcal P_s,
\qquad
\mathcal P_s
:=-\bigl\langle\Lambda^su,
\Lambda^s\bigl((u\cdot\nabla)u+\nabla p\bigr)\bigr\rangle.
\]
Repeated differentiation gives
\[
H^{(n)}
=(-2\nu)^nE_{n+\frac12}
+\sum_{j=0}^{n-1}(-2\nu)^j
\partial_t^{\,n-1-j}\mathcal P_{j+\frac12}.
\]
The \(n\)-th temporal response contains the spatial energy at height \(n+\tfrac12\), together with every differentiated nonlinear and pressure transfer. As the temporal order rises, the Sobolev height rises with it.

Keep the finite-energy row with those half-integer heights. Then
\[
L^2(\mathbb R^3)
\cap\bigcap_{n=0}^{\infty}\dot H^{n+\frac12}(\mathbb R^3)
=\bigcap_{m=0}^{\infty}H^m(\mathbb R^3)
=H^\infty(\mathbb R^3).
\]
In three dimensions,
\[
s>k+\frac32
\quad\Longrightarrow\quad
H^s(\mathbb R^3)\hookrightarrow C^k(\mathbb R^3),
\]
so \(H^\infty(\mathbb R^3)\hookrightarrow C^\infty(\mathbb R^3)\). A singularity would need the temporal ladder to run toward spatial roughness. If every finite coupled row can be kept on the whole interval, the ladder runs into smoothness instead.

Take one finite temporal order \(N\) and one finite spatial order \(m\). Differentiate only to those orders and keep the resulting velocity, nonlinear, pressure, and viscous terms in one block. The pressure-complete recurrence fixes how its entries overlap; the finite mixed estimate controls the block on the whole interval. In particular,
\[
V_N\in L^2(0,T_*;H^{m+1}),
\qquad
V_{N+1}\in L^2(0,T_*;H^{m-1}),
\]
and
\[
\sup_{0<T<T_*}\mathcal R_{N,m}(u;T)<\infty.
\]
Viscosity spans the two spatial levels in the adjacent pair; the nonlinear products and pressure remain inside the joined row. The constant may depend on \(N\) and \(m\). No common constant over the infinite tower is needed.

Repeat the finite argument for every \(N\) and \(m\). Every block comes from one velocity history generated by \(u_0\), so entries shared by two blocks already agree. They fit into
\[
u_0
\longrightarrow
\{\mathcal R_{N,m}[u_0]\}_{N,m<\infty}
\longrightarrow
\mathcal I[u_0]
:=
\bigcap_{N,m<\infty}H_t^N(0,T_*;H_x^m).
\]
Compatibility turns the finite estimates into one all-orders intersection without taking an infinite-order limit of their constants.

For each finite \(m\), the intersection contains
\[
u\in L^2(0,T_*;H^{m+1}),
\qquad
u_t\in L^2(0,T_*;H^{m-1}).
\]
The Hilbert-space trace theorem gives
\[
u\in C([0,T_*];H^m).
\]
Because these traces all belong to one velocity field,
\[
u_*:=u(T_*)\in\bigcap_{m<\infty}H^m=H^\infty.
\]
The pressure and velocity recurrences rebuild the complete endpoint jet:
\[
Q_N
=R_iR_j\sum_{a=0}^{N}\binom NaV_{a,i}V_{N-a,j},
\]
\[
V_{N+1}
=\nu\Delta V_N
-\sum_{a=0}^{N}\binom Na(V_a\cdot\nabla)V_{N-a}
-\nabla Q_N.
\]
For every \(N\), both \(V_N\) and \(Q_N\) lie in \(H^\infty\) at \(T_*\). Smooth local existence restarts the solution from \(u_*\), and uniqueness glues that restart to the original velocity history. The solution continues past \(T_*\), contradicting maximality. The finite terminal time cannot occur.

This is not an a priori estimate proof. It is an intersection-to-estimate proof. Viscosity relates the velocity's time derivatives to its spatial Sobolev derivatives. Taken through every finite order, those two directions meet in a common intersection. At that intersection, \(H^\infty(\mathbb R^3)\hookrightarrow C^\infty(\mathbb R^3)\) gives smoothness, and the familiar estimates can be taken from whichever finite Sobolev height they require. The estimates come from the intersection, not the other way around.

\end{document}
